% !TEX root = ../main.tex
\section{Background}
\label{sec:background}

The goal of attributing outcomes to system interventions in long-horizon, goal-oriented workflows also forms the foundation of controlled experiments and causal inference. The design and analysis of online A/B and A/B/n tests are well established \citep{kohavi2020trustworthy}; causal attribution is formalised in the potential-outcomes framework \citep{rubin2005causal, imbens2015causal} and in structural causal models \citep{pearl2009causality}. Key challenges include sample size and power, variance reduction, and segment-level analysis when treatment effects are small or diluted across long user journeys. In this paper we adopt a generative model (a nested Markov chain) so that outcome metrics and trajectory length have well-defined distributions; this lets us derive sample size and allocation in closed form and compare systems with confidence intervals on metric differences.

Absorbing Markov chains and the fundamental matrix provide closed-form expressions for absorption probabilities, expected hitting times, and related quantities \citep{kemeny1976finite, norris1997markov}. Phase-type distributions characterise the time until absorption in finite-state Markov chains and are widely used in applied probability \citep{neuts1981matrix}. We use the same machinery: our top-level chain has absorbing states (finished, abandoned) and outcome states (publish, subscribe), and our nested chain per goal has transient and absorbing sub-states (succeeded, failed), so that quantities such as $\mathbb{P}(\text{subscribe})$ and expected trajectory length can be derived and used for sample-size and comparison criteria.

\citet{takanobu2020goal} show that component-wise, single-turn evaluation can be inconsistent with overall goal-oriented performance and argue for system-wise evaluation and simulated users despite simulator-human discrepancy. We adopt a goal-level decomposition but model goals and nested sub-states explicitly in the Markov structure, so that segmentation is defined by the model rather than by ad hoc clustering of turns, and interventions can be applied and measured at the goal level.

To the best of our knowledge, no prior work combines a nested Markov chain over goals and per-goal sub-states with closed-form outcome attribution, sample size, optimal allocation, and leverage (sensitivities) under a finite set of user personas for agent-human workflows. In the following sections we describe the model, inference under label noise, and the resulting design levers.
